\chapter{Random Variables: The Idea}

In the previous chapters, we reasoned about \emph{events}: a medical test, a coin
flip, or a spade drawn from a deck. But events are strictly yes/no questions,
assigning each outcome a probability between $0$ and $1$.

Yet most questions in science and daily life are not yes/no. We do not just want
to know \emph{whether} a gambler wins; we want to know \emph{how much} they win. We
do not just want to know whether a clinical trial shows an effect; we want the
\emph{number} of patients who improved. We do not just want to know whether a network packet
arrives; we want the \emph{delay} in milliseconds. The world speaks to us in
\emph{numbers attached to random outcomes}, and the mathematical object that captures
``a number attached to a random outcome'' is the \textbf{random variable}.

At bottom, a random variable is simply a function defined on the sample space.
Much of the rest of the book builds on what that one idea lets us do.

\section{From outcomes to numbers}
\label{ch06:outcomes-to-numbers}

Start with a story you already understand and watch the new idea grow out of
it. Roll two fair six-sided dice. The sample space -- the set of all possible
outcomes -- is the $36$ ordered pairs
\[
  \Omega = \set{(1,1),(1,2),\dots,(6,6)},
\]
each equally likely with probability $1/36$. So far this is Chapter-3 material.

Now ask a natural question: \emph{what is the sum of the two dice?} The sum is not an
outcome in $\Omega$; the outcomes are pairs. The sum is a \emph{number we compute from}
an outcome. If the dice come up $(3,5)$, the sum is $8$. If they come up $(6,6)$, the
sum is $12$. There is a rule -- a machine -- that eats an outcome and spits out a
number.

That machine is a random variable. Call it $S$. Formally, $S$ is a function from the
sample space to the real numbers:
\[
  S : \Omega \to \R, \qquad S\big((a,b)\big) = a + b.
\]
Both words in ``random variable'' are slightly misleading: there is nothing
random or variable about the function $S$ itself. $S$ is a fixed,
deterministic rule. The randomness lives entirely in \emph{which outcome
$\omega \in \Omega$ gets fed in}. Once nature picks $\omega$, the value $S(\omega)$ is
completely determined. The randomness comes in the front door (which $\omega$?) and
the function carries it straight through to the number $S(\omega)$.

\begin{intuitionbox}
Think of a random variable as a \emph{measuring instrument} pointed at a random
experiment. The experiment produces an outcome $\omega$ -- the full, detailed state of
the world. The instrument reads off one number from that state. A thermometer in a
turbulent room, a scoreboard tracking a game, the meter on a taxi: each is a fixed
device, but because the world it measures is uncertain, the reading is uncertain too.
The random variable \emph{is} the device, not the reading.
\end{intuitionbox}

\begin{figure}[H]
\centering
\begin{tikzpicture}[>={Stealth[round]}, font=\small]
% sample space blob on the left
\node[draw, ellipse, minimum width=3.6cm, minimum height=4.2cm,
      fill=sourceblue!8, draw=sourceblue!80] (om) at (0,0) {};
\node[sourceblue!60!black] at (0,1.95) {$\Omega$};
\node (a) at (0,1.15) {$(2,6)$};
\node (b) at (0,0.5)  {$(3,5)$};
\node (c) at (0,-0.15){$(4,4)$};
\node (d) at (0,-0.8) {$(5,3)$};
\node (e) at (0,-1.45){$(6,2)$};
% number line on the right
\draw[->, thick] (5.0,-1.9) -- (5.0,2.2);
\foreach \y/\v in {-1.6/6, -0.8/7, 0/8, 0.8/9, 1.6/10}{
  \draw (4.9,\y) -- (5.1,\y) node[right] {$\v$};
}
\node[targetgreen!55!black] at (5.0,2.5) {$\R$};
% arrows: all five pairs map to 8
\foreach \n in {a,b,c,d,e}{
  \draw[->, sourceblue!75] (\n.east) to[out=0,in=180] (4.9,0);
}
\node[align=center, font=\footnotesize] at (2.5,-2.5)
  {$S\big((a,b)\big)=a+b$};
\node[sourceblue!90!black, font=\small\bfseries] at (2.5,0.95) {$S$};
\end{tikzpicture}
\caption{A random variable is a function. The instrument $S$ reads each detailed
outcome $\omega\in\Omega$ and reports a single number $S(\omega)\in\R$. All five
pairs summing to $8$ are collapsed onto the value $8$; the event $\set{S=8}$ is exactly
their pre-image (the collection of all input outcomes that map to that output value).
The randomness enters through which $\omega$ nature selects; the function itself is fixed.}
\label{ch06:fig-rv-map}
\end{figure}

\begin{definition}[Random variable, discrete first pass]
\label{ch06:def-rv}
Let $\Omega$ be the sample space of a random experiment. A \textbf{random variable} is
a function $X : \Omega \to \R$ that assigns a real number $X(\omega)$ to each outcome
$\omega \in \Omega$. We write $X$, $Y$, $S$, $N$, \dots (capital letters) for the
function itself, and lowercase $x$, $y$, $s$, $n$ for particular numerical values it
might take.
\end{definition}

The capital/lowercase convention is the single most
important piece of notation in the subject, and understanding it clearly now will
prevent confusion later. $X$ is the random thing -- the device. $x$ is a
plain number -- a possible reading. The expression $\set{X = x}$ then means something
precise: it is the \emph{event}
\[
  \set{X = x} \;=\; \set{\omega \in \Omega : X(\omega) = x},
\]
the collection of all outcomes that the device would read as $x$. For our dice,
$\set{S = 8}$ is the set of pairs summing to $8$:
\[
  \set{S = 8} = \set{(2,6),(3,5),(4,4),(5,3),(6,2)},
\]
five outcomes out of thirty-six. And here is the bridge from the old world to the new:
because $\set{S = 8}$ is just an ordinary event, we can ask for its probability using
everything we already know.
\[
  \P(S = 8) = \frac{\abs{\set{S=8}}}{\abs{\Omega}} = \frac{5}{36}.
\]

\begin{walkbox}
\textbf{Pause and think:} Roll the same two fair dice, but now define the instrument
$D\big((a,b)\big) = \abs{a - b}$, the absolute difference between the dice. What is the
event $\set{D = 0}$ as a set of outcomes, and what is its probability $\P(D = 0)$?
(Answer: $\set{D = 0} = \set{(1,1), (2,2), (3,3), (4,4), (5,5), (6,6)}$, the six doubles.
Because each of the $36$ outcomes is equally likely, $\P(D = 0) = 6/36 = 1/6$.)
\end{walkbox}

\begin{happeningbox}
A random variable is a \emph{relabeling} of the sample space. The original outcomes
might be pairs of dice, sequences of coin flips, points in a square, or whatever messy
thing the experiment naturally produces. The random variable collapses all of that
detail down onto a number line. Outcomes $(2,6)$, $(3,5)$, $(4,4)$, $(5,3)$, $(6,2)$
are five genuinely different states of the world, but $S$ glues them together: as far
as $S$ is concerned they are all just ``$8$,'' as Figure~\ref{ch06:fig-rv-map} shows.
This deliberate forgetting is where the leverage comes from. We throw away the
information we do not need and keep the number we do.
\end{happeningbox}

\subsection{Why bother with the function?}

You might reasonably object: I have been computing things like ``probability the sum
is $8$'' since grade school without ever uttering the word ``function.'' Why introduce
the formal machine $S : \Omega \to \R$ at all?

Three reasons, each of which will pay off repeatedly.

\textbf{First, it lets one experiment feed many questions.} From the \emph{same} roll
of two dice we can define $S = a+b$ (the sum), $M = \max(a,b)$ (the larger die),
$D = \abs{a-b}$ (the absolute difference), and $T = \ind\set{a = b}$ (an indicator of
doubles). These are four different functions on the \emph{one} sample space $\Omega$.
Each is a different instrument reading the same underlying randomness. Keeping the
sample space fixed and letting the functions multiply is enormously economical;
in particular, it directly allows us to ask how these readings relate to each other (Are $S$ and $M$
related? Does knowing $D$ tell us anything about $S$?), which is the entire content of
later chapters on joint distributions and correlation.

\textbf{Second, it lets us do arithmetic with randomness.} If $X$ and $Y$ are random
variables on the same $\Omega$, then $X + Y$, $XY$, $3X - 2$, $\max(X,Y)$, and
$e^{X}$ are \emph{also} random variables -- they are just new functions on $\Omega$,
built pointwise: $(X+Y)(\omega) = X(\omega) + Y(\omega)$. Treating randomness as
functions means we inherit all of algebra for free. ``The total winnings over ten
games'' becomes a sum of ten random variables, and we can manipulate it like any other
algebraic expression.

\textbf{Third, it is the abstraction that lets us forget the sample space entirely.}
This is the deepest reason. Once we know how a random variable distributes its
probability over the number line, we almost never need to look back at $\Omega$ again.
Whether ``the number of heads in $10$ flips'' came from coins, from a sequence of
yes/no survey answers, or from defective items on an assembly line, the \emph{numbers}
behave identically. The random variable is the interface; $\Omega$ is the hidden
implementation. The next section makes this concrete.

\begin{connectionbox}
If you have read the companion book on linear algebra, the move here will feel
familiar. There, a linear map is the central object and a matrix is merely one
representation of it in a chosen basis. Here, the \emph{distribution} (what comes next)
is the central object, and the sample space $\Omega$ is merely one underlying ``basis''
that realizes it. Mathematicians habitually push the concrete outcomes into the
background and reason about the structure. A random variable is the first
big instance of that habit in probability.
\end{connectionbox}

\section{The distribution of a random variable}
\label{ch06:distribution}

Now comes the conceptual leap that makes random variables indispensable. We have a
function $X : \Omega \to \R$. The function induces, on the number line, a \emph{spread
of probability}: each possible value $x$ gets the probability of the event that lands
on it. This spread is called the \textbf{distribution} of $X$. It answers, in one
package, every question of the form ``how likely is $X$ to fall here?''

\begin{definition}[Distribution]
\label{ch06:def-distribution}
The \textbf{distribution} of a random variable $X$ is the rule that assigns to each
(reasonable) set $B \subseteq \R$ the number
\[
  \P(X \in B) \;=\; \P\big(\set{\omega : X(\omega) \in B}\big).
\]
It tells us how the total probability mass of $1$ is distributed along the real line.
\end{definition}

For a discrete random variable -- one whose possible values form a finite or countable
set -- the distribution is captured completely by listing how much probability sits on
each individual value. That list is the \textbf{probability mass function}.

\begin{definition}[Probability mass function (PMF)]
\label{ch06:def-pmf}
Let $X$ be a discrete random variable. Its \textbf{probability mass function} is
\[
  p_X(x) \;=\; \P(X = x) \;=\; \P\big(\set{\omega : X(\omega) = x}\big),
  \qquad x \in \R.
\]
The PMF is positive only on the (countable) set of values $X$ actually takes, called
the \emph{support} of $X$.
\end{definition}

A PMF cannot be just any function. To be a legitimate distribution of probability it
must obey two rules, and these rules are simply the probability axioms, transported
through the function $X$.

\begin{proposition}[Valid PMF]
\label{ch06:prop-valid-pmf}
A function $p_X$ is the PMF of some discrete random variable if and only if
\begin{enumerate}
  \item $p_X(x) \ge 0$ for every $x$ (probabilities are never negative), and
  \item $\displaystyle\sum_{x} p_X(x) = 1$ (the total probability is one),
\end{enumerate}
where the sum runs over the support of $X$.
\end{proposition}

\begin{whybox}
Why must the masses sum to $1$? Because the events $\set{X = x}$, as $x$ ranges over
the support, are \emph{disjoint} (an outcome cannot make the instrument read both $3$
and $5$ at once) and \emph{exhaustive} (every outcome reads \emph{some} value). They
partition $\Omega$ -- meaning they slice the sample space into non-overlapping pieces
that together cover every possible outcome. By countable additivity -- the third
probability axiom from Chapter~3 -- the probabilities of a partition add up to
$\P(\Omega) = 1$:
\[
  \sum_x p_X(x) = \sum_x \P(X = x) = \P\!\left(\bigcup_x \set{X=x}\right) = \P(\Omega) = 1.
\]
Nonnegativity is just the first axiom passed along. So a PMF is nothing more
exotic than the probability measure on $\Omega$, repackaged on the number line. No new
axioms are needed; this is simply the original probability measure pushed forward onto
the real line.
\end{whybox}

\begin{intuitionbox}
\textbf{The physical beam: why ``mass''?}
Why do probabilists speak of probability \emph{mass} rather than just numbers?
Picture the real number line as a thin, weightless rigid beam or ruler. At each location
$x$ where $X$ can land, imagine placing a small physical weight of magnitude $p_X(x)$.
The total weight of all objects placed on the beam is exactly $\sum_x p_X(x) = 1$ unit.

The distribution describes how that unit of mass is spread out across the beam. If
probability is concentrated near a single number, the beam is heavy there; if it is
spread evenly, the weight is dispersed. When we later ask where this beam would naturally
balance on a triangular pivot, the answer will be the \emph{center of mass} -- the
expected value $\E[X]$ introduced in Chapter~8. Thinking of probabilities
as physical masses on a line grounds our spatial intuition in mechanics.
\end{intuitionbox}

\begin{figure}[H]
\centering
\begin{tikzpicture}[font=\small, >=stealth]
  % Rigid beam
  \draw[thick, fill=gray!15, rounded corners=2pt] (-0.5,0) rectangle (9.5,0.22);
  \node[anchor=east, font=\footnotesize\color{black!70}] at (-0.6,0.11) {Rigid beam};

  % Tick marks and coordinate labels for positions 1, 2, 4, 5
  \foreach \x/\val in {1/1, 3/2, 7/4, 9/5} {
    \draw[thick, gray!70] (\x, 0) -- (\x, -0.15) node[below, font=\footnotesize] {$\val$};
  }
  \node[font=\footnotesize\color{black!70}] at (10.2,-0.3) {coordinate $x$};

  % Weights placed on the beam
  % Position 1: mass 0.1
  \draw[fill=sourceblue!25, draw=sourceblue!80, rounded corners=2pt, thick] (0.6,0.22) rectangle (1.4,0.72);
  \node[font=\scriptsize\bfseries\color{sourceblue!80!black}] at (1,0.47) {$0.1$};
  \node[above, font=\footnotesize, align=center] at (1,0.8) {$p_X(1) = 0.1$};

  % Position 2: mass 0.4
  \draw[fill=sourceblue!35, draw=sourceblue!80, rounded corners=2pt, thick] (2.5,0.22) rectangle (3.5,1.52);
  \node[font=\scriptsize\bfseries\color{sourceblue!80!black}] at (3,0.87) {$0.4$};
  \node[above, font=\footnotesize, align=center] at (3,1.6) {$p_X(2) = 0.4$};

  % Position 3: mass 0.2
  \draw[fill=sourceblue!25, draw=sourceblue!80, rounded corners=2pt, thick] (4.6,0.22) rectangle (5.4,0.97);
  \node[font=\scriptsize\bfseries\color{sourceblue!80!black}] at (5,0.6) {$0.2$};
  \node[above, font=\footnotesize, align=center] at (5,1.05) {$p_X(3) = 0.2$};

  % Position 4: mass 0.0 (no block)
  \node[above, font=\footnotesize\color{gray}] at (7,0.3) {$p_X(4) = 0$};

  % Position 5: mass 0.3
  \draw[fill=sourceblue!30, draw=sourceblue!80, rounded corners=2pt, thick] (8.5,0.22) rectangle (9.5,1.32);
  \node[font=\scriptsize\bfseries\color{sourceblue!80!black}] at (9,0.77) {$0.3$};
  \node[above, font=\footnotesize, align=center] at (9,1.4) {$p_X(5) = 0.3$};

  % Fulcrum at center of mass
  % Center of mass: 1*0.1 + 2*0.4 + 3*0.2 + 4*0 + 5*0.3 = 0.1 + 0.8 + 0.6 + 0 + 1.5 = 3.0
  % x = 3 corresponds to coordinate 5.0 on canvas
  \filldraw[fill=quotientorange!30, draw=quotientorange!90, thick] (5,0) -- (4.65,-0.65) -- (5.35,-0.65) -- cycle;
  \draw[thick, quotientorange!80!black] (4.4,-0.65) -- (5.6,-0.65);
  \node[font=\footnotesize\bfseries\color{quotientorange!90!black}] at (5,-0.38) {$3$};
  \node[below=2pt, font=\footnotesize\bfseries\color{quotientorange!80!black}, align=center] at (5,-0.72)
    {Fulcrum at $x = 3$ ($\E[X] = 3.0$)\\$\sum x \cdot p_X(x) = 3.0$};

  % Total mass indicator
  \node[draw=targetgreen!80, fill=targetgreen!10, rounded corners=2pt, font=\footnotesize, align=center] at (1.2,-1.2)
    {Total mass:\\$\sum_x p_X(x) = 1.0$};
\end{tikzpicture}
\caption{Probability as physical mass on a rigid balance beam. Each outcome $x$ carries a physical weight $p_X(x)$. The total mass distributed along the beam is exactly $\sum_x p_X(x) = 1$. The beam balances at its mechanical center of mass (the fulcrum), previewing the expected value $\E[X]$ introduced in Chapter~8.}
\label{ch06:fig-balance-beam}
\end{figure}

\subsection{The dice sum, fully worked}

Let us build the entire distribution of the dice sum $S$ and look at it. We need
$\P(S = s)$ for each $s$ from $2$ through $12$. Each is a count of pairs summing to $s$,
divided by $36$.

\begin{example}[Distribution of the sum of two dice]
\label{ch06:ex-dice-sum}
Counting the pairs $(a,b)$ with $a+b = s$:
\[
\begin{array}{c|ccccccccccc}
s & 2 & 3 & 4 & 5 & 6 & 7 & 8 & 9 & 10 & 11 & 12 \\ \hline
\#\set{S=s} & 1 & 2 & 3 & 4 & 5 & 6 & 5 & 4 & 3 & 2 & 1
\end{array}
\]
so the PMF is
\[
  p_S(s) = \frac{\#\set{S=s}}{36},
  \qquad s = 2,3,\dots,12.
\]
As a sanity check, the counts add to $1+2+\cdots+6+\cdots+2+1 = 36$, so
$\sum_s p_S(s) = 36/36 = 1$. The distribution is a perfect triangle: it rises to a
peak at $s = 7$, the most likely sum, and falls off symmetrically. There is exactly
one way to get a $2$ (snake eyes) or a $12$ (boxcars), but six ways to get a $7$,
which is why $7$ rules at the craps table.
\end{example}

\begin{figure}[H]
\centering
\begin{tikzpicture}
\begin{axis}[
  width=12cm, height=6.2cm,
  title style={font=\small},
  title={Probability Mass Function of the Dice Sum $S = a + b$},
  ybar, bar width=14pt,
  xlabel={value $s$ of the sum}, ylabel={$p_S(s)$},
  xmin=1.2, xmax=12.8,
  xtick={2,3,4,5,6,7,8,9,10,11,12},
  ymin=0, ymax=0.20,
  ytick={0,0.05,0.10,0.15},
  yticklabel style={/pgf/number format/fixed, /pgf/number format/precision=2},
  ymajorgrids, grid style={gray!18},
  axis lines=left,
]
\addplot[fill=distblue!70, draw=distblue!90] coordinates {
  (2,0.02778) (3,0.05556) (4,0.08333) (5,0.11111) (6,0.13889)
  (7,0.16667) (8,0.13889) (9,0.11111) (10,0.08333) (11,0.05556) (12,0.02778)
};
\end{axis}
\end{tikzpicture}
\caption{The PMF of $S$, the sum of two fair dice. Each bar's height is the
probability mass on that value. The masses are nonnegative and add to $1$; the shape
is the famous craps triangle, peaking at $7$.}
\label{ch06:fig-dice-pmf}
\end{figure}

Figure~\ref{ch06:fig-dice-pmf} is your first \emph{distribution plot}, and it deserves
a moment. The bar chart \emph{is} the distribution -- it shows, at a glance, where the
random variable likes to land. Everything we will eventually compute about $S$ -- its
average value, its spread, the chance it exceeds $9$ -- is already encoded in these
eleven bars. The $36$-element sample space has dropped out of sight entirely; the
triangle is all we need.

\begin{walkbox}
Use Figure~\ref{ch06:fig-dice-pmf} to read off $\P(S \ge 10)$ without any new
computation. (Add the masses on $10$, $11$, $12$.) Then find $\P(S \text{ is even})$.
Which is bigger, $\P(S = 7)$ or $\P(S \in \set{2,3,11,12})$? Answers:
$\P(S \ge 10) = (3+2+1)/36 = 1/6$; the even sums $2,4,6,8,10,12$ carry
$(1+3+5+5+3+1)/36 = 18/36 = 1/2$; and the four extreme values carry
$(1+2+2+1)/36 = 6/36 = 1/6 = \P(S=7)$ exactly -- a tie.
\end{walkbox}

\subsection{Same distribution, different worlds}

Here is the payoff promised earlier. Consider three completely different experiments.

\begin{example}[Three experiments, one distribution]
\label{ch06:ex-same-dist}
\textbf{(A)} Flip a fair coin three times; let $X$ be the number of heads.
\textbf{(B)} A factory line produces three independent items, each defective with
probability $1/2$; let $Y$ be the number of defectives.
\textbf{(C)} Three independent voters each answer ``yes'' with probability $1/2$; let
$Z$ be the number of yeses.

The sample spaces are different in flavor -- coins, widgets, ballots -- but each
$X,Y,Z$ counts ``successes'' in three independent fair trials. For $X$ (and identically
for $Y,Z$):
\[
  p_X(0) = \tfrac{1}{8},\quad p_X(1) = \tfrac{3}{8},\quad
  p_X(2) = \tfrac{3}{8},\quad p_X(3) = \tfrac{1}{8}.
\]
(There are $2^3 = 8$ equally likely outcomes; $\binom{3}{k}$ of them have $k$
successes.) All three random variables have \emph{exactly the same PMF}. They live on
different sample spaces, but their distributions are identical.
\end{example}

This is why probabilists give distributions their own names. We say $X$, $Y$, and $Z$
are each \emph{Binomial} with parameters $n = 3$ and $p = 1/2$, written
$X \sim \Bin(3, 1/2)$. The symbol $\sim$ reads ``is distributed as.'' From this point
forward, the moment you hear ``count of successes in $n$ independent trials, each
succeeding with probability $p$,'' you may forget the coins and widgets and ballots
entirely and reach for the $\Bin(n,p)$ distribution. The named distribution is a reusable
template; the next chapter is devoted to building a small library of such templates.

\begin{aestheticbox}
This is the moment probability becomes a science rather than a bag of puzzles. By
focusing on the \emph{distribution} -- the spread of probability on the number line --
we discover that wildly different situations share the same mathematical heart. The
same bell curve governs measurement errors in astronomy, heights in a population, and
the noise in a financial return. The same Poisson law counts radioactive decays,
typos per page, and customers per minute. Recognizing the distribution beneath the
story is much of what a working scientist actually does, and the random-variable
abstraction is what makes that recognition possible.
\end{aestheticbox}

\begin{historybox}
The concept of a ``random variable'' is surprisingly modern. For nearly three centuries
after Blaise Pascal and Pierre de Fermat founded mathematical probability in their 1654
letters, nobody had a formal name for the numerical reading produced by an experiment.
Pioneers like Gerolamo Cardano, Christiaan Huygens, Jakob Bernoulli, Abraham de Moivre,
and Pierre-Simon Laplace computed fair stakes, game durations, and probabilities of
winning with remarkable ingenuity, but their tools were strictly combinatorics and
calculus tied to physical objects -- dice faces, playing cards, and colored balls in
urns. They were stuck because every new problem forced them to reason about raw physical
outcomes from scratch. Without an abstract bridge separating the physical trial
$\Omega$ from the measured quantity $X$, they could not easily add uncertain numbers
($X+Y$), study transformations, or state universal limit laws without drowning in the
mechanical details of specific games.

In the late nineteenth century, the Russian school -- Pafnuty Chebyshev and his brilliant
students Andrei Markov and Aleksandr Lyapunov -- took a decisive step forward. They
isolated ``random quantities'' ($\text{slu}\check{c}\text{ajnaya velichina}$) as
mathematical objects in their own right, proving the first general central limit
theorems about sums of independent variables.

Yet a sharp controversy soon erupted over the very legitimacy of the subject. In 1900,
David Hilbert presented his famous Sixth Problem, challenging mathematicians to put
probability on the same rigorous axiomatic foundation as geometry. For decades, the
field remained contested. The applied mathematician Richard von Mises argued in 1919
that probability must be defined through long-run frequencies in infinite trial
sequences, which he called ``collectives.'' But von Mises's framework sparked bitter
disputes: how could one define an infinite sequence as genuinely ``random'' without
falling into circular logic? Pure mathematicians often dismissed probability as mere
gambling heuristics or physics rather than true mathematics.

The impasse was finally broken in 1933 by the 30-year-old Andrey Kolmogorov in his
monograph \emph{Grundbegriffe der Wahrscheinlichkeitsrechnung} (Foundations of the Theory
of Probability). Kolmogorov realized that probability could be grounded entirely in
measure theory -- the mathematical framework developed by Emile Borel and Henri Lebesgue
to measure lengths, areas, and volumes. Under Kolmogorov's axioms, the sample space is
just an abstract set $\Omega$, events are measurable subsets, and a random variable is
simply a \emph{measurable function} -- a rule assigning a real number to each outcome
that is well-behaved enough to let us measure the probability of its values.

This simple definition detached probability from coins and urns forever. It connected probability to
measure-theoretic integration and functional analysis, laying the foundation for stochastic
calculus (the mathematics of random motion that powers quantitative finance and
diffusion models), martingale theory, and rigorous statistical learning.

Even today, the boundaries of the concept remain active frontiers. In quantum
physics, physical quantities like position and momentum cannot be measured simultaneously
because the order of measurement matters ($XY \ne YX$). This non-commutativity led John
von Neumann and others to develop ``non-commutative probability theory,'' where random
variables become linear operators on infinite-dimensional vector spaces rather than
ordinary functions on a sample space. Meanwhile, in modern cosmology and algorithmic
information theory, philosophers and mathematicians still debate what it means to assign
a probability distribution to unique, unrepeatable systems -- such as the initial state
of the universe itself. What began as a dice player's tally took three centuries to
clarify, and its horizons are still expanding.
\end{historybox}

\section{A first look at the CDF}
\label{ch06:cdf}

The PMF answers ``how much mass sits exactly at $x$?'' Often the more useful question
is ``how much mass sits at or below $x$?'' -- the running total. That accumulated
probability is the \textbf{cumulative distribution function}, or CDF, and it is the one
description of a distribution that works for \emph{every} random variable, discrete or
continuous, without modification. We will lean on it heavily once continuous variables
arrive, so meet it now in the friendly discrete setting.

\begin{definition}[Cumulative distribution function (CDF)]
\label{ch06:def-cdf}
The \textbf{CDF} of any random variable $X$ is
\[
  F_X(x) \;=\; \P(X \le x), \qquad x \in \R.
\]
For a discrete $X$ it is a running sum of the PMF:
\[
  F_X(x) \;=\; \sum_{t \le x} p_X(t).
\]
\end{definition}

For the dice sum $S$, the CDF climbs in jumps. Below $2$ it is flat at $0$; at $s=2$ it
jumps up by $p_S(2) = 1/36$; at $s=3$ it jumps by $2/36$; and so on, reaching $1$ at
$s=12$ and staying there. The \emph{size of each jump is exactly the PMF mass at that
point} -- so you can read the entire PMF off the jumps of the CDF, and vice versa. They are two views of
the same distribution.

\begin{figure}[H]
\centering
\begin{tikzpicture}
\begin{axis}[
  width=12cm, height=6cm,
  title style={font=\small},
  title={Cumulative Distribution Function of the Dice Sum $F_S(x) = \P(S \le x)$},
  xlabel={$x$}, ylabel={$F_S(x) = \P(S \le x)$},
  xmin=1, xmax=13, ymin=0, ymax=1.05,
  xtick={2,3,4,5,6,7,8,9,10,11,12},
  ytick={0,0.25,0.5,0.75,1},
  axis lines=left,
  legend pos=south east,
]
% staircase: horizontal treads (closed left, value attained at the integer)
\addplot[distblue, very thick] coordinates {(1,0) (2,0)};
\addplot[distblue, very thick] coordinates {(2,0.02778) (3,0.02778)};
\addplot[distblue, very thick] coordinates {(3,0.08333) (4,0.08333)};
\addplot[distblue, very thick] coordinates {(4,0.16667) (5,0.16667)};
\addplot[distblue, very thick] coordinates {(5,0.27778) (6,0.27778)};
\addplot[distblue, very thick] coordinates {(6,0.41667) (7,0.41667)};
\addplot[distblue, very thick] coordinates {(7,0.58333) (8,0.58333)};
\addplot[distblue, very thick] coordinates {(8,0.72222) (9,0.72222)};
\addplot[distblue, very thick] coordinates {(9,0.83333) (10,0.83333)};
\addplot[distblue, very thick] coordinates {(10,0.91667) (11,0.91667)};
\addplot[distblue, very thick] coordinates {(11,0.97222) (12,0.97222)};
\addplot[distblue, very thick] coordinates {(12,1) (13,1)};
% solid dots at the attained values (left endpoints of each new tread)
\addplot[only marks, mark=*, distblue, mark size=1.6pt] coordinates {
  (2,0.02778) (3,0.08333) (4,0.16667) (5,0.27778) (6,0.41667) (7,0.58333)
  (8,0.72222) (9,0.83333) (10,0.91667) (11,0.97222) (12,1)
};
\end{axis}
\end{tikzpicture}
\caption{The CDF of the dice sum $S$. It is a right-continuous staircase: flat between
the integers, jumping up at each value of the support. The height of the jump at $s$
equals the PMF mass $p_S(s)$. The function starts at $0$ on the far left and reaches
$1$ on the far right.}
\label{ch06:fig-dice-cdf}
\end{figure}

Three properties leap out of Figure~\ref{ch06:fig-dice-cdf}, and they hold for the CDF
of \emph{any} random variable, not just this one. We state them now and prove them
carefully; they will be old friends by the end of the book.

\begin{proposition}[Defining properties of a CDF]
\label{ch06:prop-cdf}
Every CDF $F_X$ satisfies:
\begin{enumerate}
  \item \textbf{Monotone:} if $a \le b$ then $F_X(a) \le F_X(b)$.
  \item \textbf{Limits:} $F_X(x) \to 0$ as $x \to -\infty$ and $F_X(x) \to 1$ as
        $x \to +\infty$.
  \item \textbf{Right-continuous:} $F_X(x) = \lim_{t \downarrow x} F_X(t)$ for all $x$.
\end{enumerate}
\end{proposition}

\begin{whybox}
\emph{Monotonicity} is just monotonicity of probability. If $a \le b$ then the event
$\set{X \le a}$ is a subset of $\set{X \le b}$ (anything at most $a$ is certainly at
most $b$), and probability respects containment, so
$F_X(a) = \P(X \le a) \le \P(X \le b) = F_X(b)$. The staircase can only go up.

\emph{The limits} say that pushing the threshold to $-\infty$ excludes everything
(probability $\to 0$) and pushing it to $+\infty$ includes everything (probability
$\to 1$). Formally these follow from continuity of probability -- the property from
Chapter~3 that the probability of nested shrinking events converges to the probability
of their intersection, and expanding events to their union -- applied to
$\set{X \le -n}\downarrow \varnothing$ and $\set{X \le n}\uparrow \Omega$.

\emph{Right-continuity} is the subtle one, and it is the reason we define the CDF with
$\le$ rather than $<$. Because $\set{X \le x}$ includes the value $x$ itself, the
function attains its new, higher level \emph{exactly at} the jump point -- the solid
dots in the figure sit on top of each riser. As we approach $x$ from the right, the
shrinking events $\set{X \le x + 1/n}$ decrease down to $\set{X \le x}$, so by
continuity of probability the values come back down to $F_X(x)$. The jump in $F_X$ at a
point $x$ is precisely $\P(X = x)$, which recovers the PMF as the size of the jump.
\end{whybox}

\begin{trapbox}
A maddeningly common error is sloppiness about $<$ versus $\le$ at the jump points.
For the dice sum, $\P(S \le 7) = F_S(7) = 21/36$, but $\P(S < 7) = F_S(6) = 15/36$;
the difference is exactly $p_S(7) = 6/36$, the mass that sits \emph{at} $7$. For a
\emph{discrete} variable, where to land matters because individual points carry mass.
(For a continuous variable in later chapters, single points carry zero mass and the
distinction evaporates -- but never assume that before you know your variable is
continuous.) When in doubt, draw the staircase and ask whether the dot at the jump is
included.
\end{trapbox}

The CDF also makes interval probabilities a one-line subtraction, which is its great
practical convenience. For any $a < b$,
\[
  \P(a < X \le b) = F_X(b) - F_X(a).
\]
This holds because $\set{X \le b}$ splits into the disjoint pieces $\set{X \le a}$ and
$\set{a < X \le b}$, and probabilities of disjoint events add. For the dice,
$\P(4 < S \le 9) = F_S(9) - F_S(4) = \tfrac{30}{36} - \tfrac{6}{36} = \tfrac{24}{36}
= \tfrac{2}{3}$.

\begin{walkbox}
\textbf{Pause and think:} Check the endpoints in $\P(a < X \le b) = F_X(b) - F_X(a)$.
Notice the left endpoint is \emph{strict} ($<$) while the right is \emph{non-strict} ($\le$).
How would you compute the closed interval probability $\P(a \le X \le b)$ using CDF values and point masses?
(Answer: Because the boundary point $a$ is now included, you add its individual mass back in:
$\P(a \le X \le b) = F_X(b) - F_X(a) + p_X(a)$. When $X$ takes only integer values, this is also
simply $F_X(b) - F_X(a - 1)$.)
\end{walkbox}

\section{Indicator random variables: small but mighty}
\label{ch06:indicators}

Of all random variables, the simplest is the one that can only be $0$ or $1$. It
attaches the number $1$ to outcomes in some event $A$ and the number $0$ to everything
else. It represents a yes/no question, and despite -- in
fact because of -- its simplicity it does an astonishing amount of work throughout
probability.

\begin{definition}[Indicator random variable]
\label{ch06:def-indicator}
For an event $A \subseteq \Omega$, the \textbf{indicator} of $A$ is the random variable
\[
  \ind_A(\omega) =
  \begin{cases}
    1 & \text{if } \omega \in A, \\
    0 & \text{if } \omega \notin A.
  \end{cases}
\]
Its distribution is concentrated on two points: $\P(\ind_A = 1) = \P(A)$ and
$\P(\ind_A = 0) = 1 - \P(A)$. A random variable with this two-point distribution is
called \emph{Bernoulli} with parameter $p = \P(A)$, written $\ind_A \sim \Bern(\P(A))$.
\end{definition}

\begin{figure}[H]
\centering
\begin{tikzpicture}[>={Stealth[round]}, font=\small]
  % Sample space Omega
  \draw[thick, fill=gray!8, draw=gray!60] (0,0) ellipse (2.6cm and 1.7cm);
  \node[gray!70!black, font=\bfseries] at (-1.6,1.0) {$\Omega$};

  % Event A inside Omega
  \draw[thick, fill=sourceblue!20, draw=sourceblue!80] (0.4,-0.1) ellipse (1.4cm and 1.0cm);
  \node[sourceblue!80!black, font=\bfseries] at (0.4,0.4) {$A$};
  \node[font=\footnotesize\color{sourceblue!80!black}] at (0.4,-0.1) {$\P(A) = p$};

  % Complement region label
  \node[font=\footnotesize\color{gray!60!black}] at (-1.2,-0.4) {$A^c$};
  \node[font=\tiny\color{gray!60!black}] at (-1.2,-0.7) {$\P(A^c) = 1-p$};

  % Sample points in A and A^c
  \fill[sourceblue!80!black] (1.0,-0.35) circle (1.8pt) node[below right, font=\tiny] {$\omega_1$};
  \fill[gray!60!black] (-1.6,0.3) circle (1.8pt) node[above left, font=\tiny] {$\omega_2$};

  % Real number line on the right
  \draw[->, thick] (5.8,-1.5) -- (5.8,1.8);
  \node[targetgreen!55!black, right] at (5.8,1.8) {$\R$};

  % Ticks and values at 0 and 1
  \draw[thick] (5.65,-0.6) -- (5.8,-0.6);
  \draw[thick] (5.65,0.9) -- (5.8,0.9);
  \node[above right=1pt, font=\small] at (5.8,0.9) {$1$};
  \node[above right=1pt, font=\small] at (5.8,-0.6) {$0$};

  % Arrows mapping Omega to {0,1}
  \draw[->, thick, sourceblue!80] (1.0,-0.35) to[out=15,in=180] (5.65,0.9);
  \draw[->, thick, gray!60] (-1.6,0.3) to[out=20,in=180] (5.65,-0.6);

  % Indicator function labels on arrows
  \node[sourceblue!80!black, font=\footnotesize\bfseries, fill=white, inner sep=1.5pt] at (3.5,1.05) {$\ind_A(\omega) = 1$};
  \node[gray!70!black, font=\footnotesize\bfseries, fill=white, inner sep=1.5pt] at (3.5,-0.55) {$\ind_A(\omega) = 0$};

  % Probability mass spikes on the line
  \draw[ultra thick, sourceblue!85] (5.8,0.9) -- (7.2,0.9);
  \node[right, font=\footnotesize\color{sourceblue!80!black}] at (7.2,0.9) {mass $p$};
  \draw[ultra thick, gray!65] (5.8,-0.6) -- (7.2,-0.6);
  \node[right, font=\footnotesize\color{gray!70!black}] at (7.2,-0.6) {mass $1-p$};
\end{tikzpicture}
\caption{The indicator $\ind_A$ as a measuring instrument. It partitions the sample
space into the event $A$ (carrying probability $p$) and its complement $A^c$ (carrying
probability $1-p$). The rule collapses all of $A$ onto the number $1$ and all of $A^c$
onto the number $0$, converting a qualitative event into a quantitative Bernoulli random variable.}
\label{ch06:fig-indicator-map}
\end{figure}

Why does such a trivial object matter? Because indicators convert \emph{logic} into
\emph{arithmetic}. The set operations on events -- and, or, not -- become ordinary
algebra on their indicators, and ``counting how many things happen'' becomes
``adding up indicators.'' Watch.

\begin{proposition}[The indicator dictionary]
\label{ch06:prop-indicator-dict}
For events $A, B \subseteq \Omega$, as functions on $\Omega$:
\[
  \ind_{A^c} = 1 - \ind_A, \qquad
  \ind_{A \cap B} = \ind_A \,\ind_B, \qquad
  \ind_{A \cup B} = \ind_A + \ind_B - \ind_A \ind_B.
\]
\end{proposition}

\begin{whybox}
Check each pointwise -- that is, plug in an arbitrary $\omega$ and verify both sides
agree. The complement: $\ind_{A^c}(\omega) = 1$ exactly when $\omega \notin A$, i.e.\
when $\ind_A(\omega) = 0$, which is exactly when $1 - \ind_A(\omega) = 1$. The
intersection: $\ind_A(\omega)\ind_B(\omega)$ is a product of two numbers in
$\set{0,1}$, so it equals $1$ if and only if \emph{both} factors are $1$, i.e.\ $\omega$
is in both $A$ and $B$ -- exactly the membership rule for $A \cap B$. The union: by
inclusion--exclusion at the level of single outcomes, $\omega$ is in $A \cup B$ iff it
is in $A$ \emph{or} $B$; the formula $\ind_A + \ind_B - \ind_A\ind_B$ counts $A$, adds
$B$, and subtracts the double-counted overlap, landing on $1$ precisely when at least
one holds. Logic in, algebra out.
\end{whybox}

\begin{walkbox}
\textbf{Pause and think:} Suppose two events $A$ and $B$ are mutually exclusive (disjoint),
meaning they cannot happen simultaneously ($A \cap B = \varnothing$). What does the
product $\ind_A \ind_B$ simplify to, and what happens to the union formula $\ind_{A \cup B}$?
(Answer: Because no outcome can lie in both sets, at least one indicator is always $0$,
making $\ind_A \ind_B = 0$ everywhere. The union formula then simplifies to
$\ind_{A \cup B} = \ind_A + \ind_B$ -- the algebraic counterpart to the additivity axiom
$\P(A \cup B) = \P(A) + \P(B)$.)
\end{walkbox}

A fundamental application of indicators is the \textbf{counting trick}: if you
want to count how many of several events occur, write the count as a sum of indicators.

\begin{example}[Counting by indicators]
\label{ch06:ex-counting-indicators}
Suppose we run $n$ independent trials and $A_i$ is the event ``trial $i$ is a
success.'' Let $N$ be the total number of successes. Then
\[
  N = \ind_{A_1} + \ind_{A_2} + \cdots + \ind_{A_n} = \sum_{i=1}^n \ind_{A_i}.
\]
At any outcome $\omega$, the right-hand side adds a $1$ for each successful trial and a
$0$ for each failure, totaling precisely the number of successes -- which is what $N$
means. So a complicated counting random variable has been disassembled into $n$ tiny,
identical, easy-to-understand pieces.
\end{example}

This decomposition looks almost too cheap to be useful, but it is one of the most
productive techniques in probability. The reason -- which we will develop
fully when expectation arrives in Chapter~8 -- is that the \emph{average} of a
sum is always the sum of the averages, even when the pieces are tangled together in
complicated ways. Counting problems that are difficult by brute force become straightforward
once you write the count as $\sum \ind_{A_i}$ and average term by term. Keep this in
your back pocket; we will return to it again and again.

\begin{applicationbox}
Indicators are the bridge from probability to data. In a spreadsheet of survey
responses, a column ``owns a car (yes/no)'' is literally a column of indicator values.
The \emph{fraction} of $1$s in that column -- the sample mean of the indicator -- is
your estimate of the underlying probability $\P(\text{owns a car})$. In machine
learning, the $0/1$ \emph{loss} of a classifier is the indicator of the event ``the
prediction is wrong,'' and the model's error rate is the average of that indicator
over the test set. One-hot encoding -- the standard way to feed a categorical variable
(``color $\in$ \{red, green, blue\}'') into a model -- replaces the category with a
little vector of indicators, one per possible value. Whenever you see a column of $0$s
and $1$s, you are looking at an indicator random variable, and its average is a
probability estimate.
\end{applicationbox}

\section{Functions of a random variable}
\label{ch06:functions-of-rv}

We noted earlier that $X + Y$, $X^2$, and $e^X$ are again random variables. Now we make
that precise and work out how a transformation reshapes a distribution -- a skill
we will need constantly.

If $X : \Omega \to \R$ is a random variable and $g : \R \to \R$ is an ordinary
function, then the composition $Y = g(X)$ is the random variable defined by
\[
  Y(\omega) = g\big(X(\omega)\big).
\]
First the experiment produces $\omega$, then $X$ reads off a number, then $g$ processes
that number. It is two machines wired in series, and the result is one new instrument.
``Apply $g$'' is just feeding the reading through a second device.

The central question is: given the distribution of $X$, what is the distribution of
$Y = g(X)$? For discrete variables there is a clean recipe -- \emph{gather the masses}.
The value $Y = y$ happens exactly when $X$ lands on some value that $g$ sends to $y$.
So we sum the PMF of $X$ over all such pre-images.

\begin{proposition}[PMF of a function of a discrete variable]
\label{ch06:prop-pmf-of-function}
If $X$ is discrete with PMF $p_X$ and $Y = g(X)$, then
\[
  p_Y(y) \;=\; \P(Y = y) \;=\; \sum_{x \,:\, g(x) = y} p_X(x).
\]
That is: to find the mass $Y$ puts on $y$, find every $x$ with $g(x) = y$ and add up
the mass $X$ already put on those $x$.
\end{proposition}

\begin{whybox}
The event $\set{Y = y}$ is, unwound through the definition,
$\set{\omega : g(X(\omega)) = y}$, which is the union over all $x$ with $g(x) = y$ of
the events $\set{X = x}$. These events are disjoint (different $x$'s), so their
probabilities add -- giving the sum over the pre-image $g^{-1}(\set{y})$. No mass is
created or destroyed; it is merely relabeled by $g$ and then pooled. When $g$ is
one-to-one, each $y$ has a single pre-image and the sum has one term:
$p_Y(y) = p_X(g^{-1}(y))$. When $g$ collapses several $x$'s onto one $y$, those masses
\emph{merge}.
\end{whybox}

\begin{example}[A transformation that merges mass]
\label{ch06:ex-square}
Let $X$ be uniform on $\set{-2,-1,0,1,2}$, so $p_X(x) = 1/5$ for each. Let
$Y = X^2$. The map $g(x) = x^2$ is not one-to-one: it sends both $-2$ and $2$ to $4$,
and both $-1$ and $1$ to $1$. Gathering masses:
\[
  p_Y(0) = p_X(0) = \tfrac{1}{5}, \qquad
  p_Y(1) = p_X(-1) + p_X(1) = \tfrac{2}{5}, \qquad
  p_Y(4) = p_X(-2) + p_X(2) = \tfrac{2}{5}.
\]
The support shrank from five values to three, and the masses on the merged values
doubled. Total is still $\tfrac{1}{5} + \tfrac{2}{5} + \tfrac{2}{5} = 1$, as it must
be. Squaring folded the symmetric distribution in half.
\end{example}

\begin{figure}[H]
\centering
\begin{tikzpicture}
\begin{axis}[
  width=6.2cm, height=5cm,
  ybar, bar width=12pt,
  title style={font=\small},
  title={$X \sim \Unif\set{-2,\dots,2}$},
  xlabel={$x$}, ylabel={$p_X(x)$},
  xmin=-2.7, xmax=2.7,
  xtick={-2,-1,0,1,2}, ymin=0, ymax=0.5, ytick={0,0.2,0.4},
  axis lines=left,
]
\addplot[fill=distblue!60, draw=distblue] coordinates {(-2,0.2)(-1,0.2)(0,0.2)(1,0.2)(2,0.2)};
\end{axis}
\begin{axis}[
  at={(7.2cm,0)},
  width=6.2cm, height=5cm,
  ybar, bar width=14pt,
  title style={font=\small},
  title={$Y = X^2$ (pooled masses)},
  xlabel={$y$}, ylabel={$p_Y(y)$},
  xmin=-0.6, xmax=4.6,
  xtick={0,1,4}, ymin=0, ymax=0.5, ytick={0,0.2,0.4},
  axis lines=left,
]
\addplot[fill=densityred!65, draw=densityred] coordinates {(0,0.2)(1,0.4)(4,0.4)};
\end{axis}
\end{tikzpicture}
\caption{A non-invertible transformation pools probability. On the left, $X$ is
uniform on five points. Squaring (right) merges the symmetric pairs $\pm1$ and $\pm2$,
doubling their masses and collapsing the support to $\set{0,1,4}$.}
\label{ch06:fig-square}
\end{figure}

\begin{trapbox}
The number-one beginner mistake with transformations is to assume
$p_Y(y) = p_X(g^{-1}(y))$ in general -- ``just substitute.'' That is only valid when
$g$ is one-to-one on the support. When $g$ merges values (like squaring, or taking
absolute value, or rounding), you \emph{must} sum over all pre-images. Forgetting the
second term turned $p_Y(4) = 2/5$ into a wrong $1/5$ in the example above, and the
masses would no longer sum to $1$ -- always a red flag that a transformation was done
carelessly. The summing-to-one check is your friend.
\end{trapbox}

\begin{walkbox}
\textbf{Pause and think:} Suppose $X$ is again uniform on $\set{-2,-1,0,1,2}$, but we now
apply the cubing function $Z = X^3$. Do any probability masses merge? What is the PMF of $Z$?
(Answer: No masses merge! The function $g(x) = x^3$ is strictly increasing, hence one-to-one
on the reals. Each input maps to a distinct cube: $(-2)^3 = -8$, $(-1)^3 = -1$, $0^3 = 0$,
$1^3 = 1$, $2^3 = 8$. The support of $Z$ has five points, each retaining its original
mass $p_Z(z) = 1/5$. Only non-invertible transformations collapse and pool probability.)
\end{walkbox}

A special, important case is the \textbf{affine} (linear-plus-shift) transformation
$Y = aX + b$ with $a \ne 0$. Here $g(x) = ax + b$ is one-to-one, so no merging occurs;
the distribution is merely relocated and rescaled. Each value $x$ moves to $ax + b$,
carrying its mass along untouched:
\[
  p_Y(y) = p_X\!\left(\frac{y - b}{a}\right).
\]
Geometrically, multiplying by $a$ stretches the number line (and flips it if $a < 0$),
while adding $b$ slides it. The \emph{shape} of the distribution is preserved; only its
location and scale change. This is exactly what happens when we convert Celsius to
Fahrenheit, dollars to euros, or -- a move we will make constantly -- when we
\emph{standardize} a variable by subtracting its center and dividing by its spread.

\begin{connectionbox}
Affine transformations of random variables are the probabilistic shadow of the affine
maps you met in linear algebra: a scaling (the $a$) composed with a translation (the
$b$). When we reach the normal distribution, the entire family of bell curves will turn
out to be the affine images of a single ``standard'' bell curve $\Normal(0,1)$ -- one
shape, infinitely relocated and rescaled. The habit of asking ``what does this
transformation do to the distribution?'' directly parallels the geometric question ``how
does this map move space around?'' For continuous variables the same question requires
a derivative correction (the change-of-variables / Jacobian factor), which the later
chapter on transformations develops in full. The discrete mass-gathering recipe here
is the clean prototype of that idea.
\end{connectionbox}

\section{A subtle point: which functions are allowed?}
\label{ch06:measurability}

We have been saying ``$X$ is any function $\Omega \to \R$.'' For finite or countable
sample spaces -- coins, dice, cards, counts -- that is literally true and you may take
it at face value; nothing can go wrong. But for the uncountable sample spaces that
appear once we model continuous quantities (a uniformly random point in $[0,1]$, say),
a fully arbitrary function can be pathological, and we need a mild restriction. This
subsection is a signpost, not a detour; you can skim it and return after meeting
continuous variables.

The whole point of a random variable was to let us ask ``$\P(X \in B)$.'' For that
question to have an answer, the set $\set{\omega : X(\omega) \in B}$ must be an
\emph{event} -- something our probability measure actually knows how to weigh. A
function for which this holds (for every reasonable target set $B$) is called
\textbf{measurable}, and the official definition of a random variable includes this
word.

\begin{formalbox}
Let $(\Omega, \mathcal{F}, \P)$ be a probability space, where $\mathcal{F}$ is the
collection of events the measure $\P$ is defined on. A function $X : \Omega \to \R$ is a
\textbf{random variable} if it is \emph{measurable}: for every $x \in \R$,
\[
  \set{\omega : X(\omega) \le x} \in \mathcal{F}.
\]
Equivalently, $X^{-1}(B) \in \mathcal{F}$ for every Borel set $B \subseteq \R$
(any subset formed from intervals by standard unions, intersections, and complements --
in practice, every subset of the real line you could ever encounter in applied work).
The condition guarantees that $\P(X \le x)$ -- and hence $\P(X \in B)$ for all the sets we
care about -- is well defined, which is precisely what makes the CDF and the
distribution exist.
\end{formalbox}

Do not let this worry you. Every function you will ever write down by an explicit
formula -- sums, products, maxima, polynomials, exponentials, indicators of sensible
events, limits of these -- is measurable. Measurability is a safety rail that rules out
artificially constructed monsters, not a constraint on real modeling. We mention it
only so that the terminology is familiar when you encounter it later, and so you understand
\emph{why} the careful definition exists: it is the price of guaranteeing that the question
``$\P(X \in B)$?'' always has an answer. The full machinery is the subject of the
measure-theoretic chapters at the end of the book; until then, ``a random variable is a
number-valued reading of a random outcome'' will serve you perfectly.

\begin{paradoxbox}
\textbf{The two-envelope puzzle, and why ``a random variable'' is not enough.}
Here is a puzzle that quietly abuses the random-variable idea. Two envelopes contain
money; one holds twice as much as the other, but you do not know which. You pick one
and find \$$100$. Should you switch? A tempting argument: the other envelope holds
either \$$50$ or \$$200$, each ``equally likely,'' so its expected contents are
$\tfrac12(50) + \tfrac12(200) = \$125 > \$100$ -- always switch! But by symmetry the
same argument applies \emph{before} you open anything, suggesting you should switch
endlessly, which is absurd.

The resolution lives exactly in the ideas of this chapter. The flawed step assumes
a \emph{distribution} for the amount of money that does not actually exist. To say ``the
other envelope is \$$50$ or \$$200$ with probability $1/2$ each'' you must specify a
genuine random variable -- a real distribution over how much money was put in the
envelopes in the first place. The instant you write down any \emph{valid} PMF for the
smaller amount (one whose masses are nonnegative and sum to $1$, as
Proposition~\ref{ch06:prop-valid-pmf} demands), the apparent ``always \$$125$''
symmetry breaks: for some observed values switching helps and for others it hurts, and
they balance. No legitimate distribution makes switching always better. The paradox is
not about envelopes; it is about pretending a random variable exists when no valid
distribution backs it up. \emph{Always ask whether your masses sum to one.}
\end{paradoxbox}

\section{Putting it together}
\label{ch06:synthesis}

One example can exercise every idea in the chapter at once: a random variable as a
function, its PMF, its CDF, an indicator decomposition, and a transformation. Here it
is.

\begin{example}[Coupons in a tiny collection]
\label{ch06:ex-coupons}
A cereal company hides one of two equally likely toys, $A$ or $B$, in each box. You buy
three boxes. Model the outcome as a string like $ABA$; there are $2^3 = 8$ equally
likely strings. Let $X$ be the number of \emph{distinct} toys you collect (so
$X \in \set{1,2}$).

\textbf{As a function.} $X(\omega) = 1$ if all three boxes match ($AAA$ or $BBB$) and
$X(\omega) = 2$ otherwise. The event $\set{X = 1}$ contains $2$ of the $8$ strings.

\textbf{PMF.} $p_X(1) = 2/8 = 1/4$ and $p_X(2) = 6/8 = 3/4$. Masses are nonnegative and
sum to $1$. (So in fact $X \sim \Bern(3/4) + 1$, a shifted coin -- or just a two-point
distribution.)

\textbf{CDF.} $F_X(x) = 0$ for $x < 1$; $F_X(x) = 1/4$ for $1 \le x < 2$; and
$F_X(x) = 1$ for $x \ge 2$. A two-step staircase, jumping by $p_X(1)$ then by $p_X(2)$.

\textbf{Indicator view.} Let $C$ be the event ``you collected both toys.'' Then
$\ind_C = \ind\set{X = 2}$, a Bernoulli($3/4$) variable, and $X = 1 + \ind_C$. We have
written the count as ``one (you always get at least one toy) plus an indicator for
having gotten the second.'' This is the counting trick in miniature.

\textbf{Transformation.} Suppose completing the set (getting both toys) wins a \$$10$
prize and otherwise you win nothing: $W = 10\,\ind_C$. This is a function of $X$ (indeed
$W = g(X)$ with $g(1) = 0, g(2) = 10$). Gathering masses, $W$ has PMF
$p_W(0) = 1/4$, $p_W(10) = 3/4$. The whole problem -- function, distribution, running
total, indicator, payoff -- has been narrated in the language of this chapter without a
single new formula.
\end{example}

\begin{recoverybox}
Take a breath and assemble the picture. A \emph{random variable} is a function from
outcomes to numbers -- a measuring device. Its \emph{distribution} is how it spreads
probability along the number line, summarized for discrete variables by the \emph{PMF}
(mass at each point, nonnegative, summing to one) and for any variable by the
\emph{CDF} (running total $\P(X \le x)$, a monotone right-continuous climb from $0$ to
$1$, whose jumps recover the PMF). \emph{Indicators} are the $0/1$ atoms that turn logic
into arithmetic and counting into summing. \emph{Functions} of a random variable are
new devices wired in series, whose distribution you find by gathering masses over
pre-images. Everything that follows -- expectation, variance, the named distributions,
the limit theorems -- is built on these five ideas. In the next chapter, we take this
measuring device and catalog the classic discrete distributions that nature produces again
and again: the Bernoulli coin flip, the Binomial count of successes, the Geometric waiting
time, and the Poisson law of rare events.
\end{recoverybox}

\section{Exercises}
\label{ch06:exercises}

\begin{exercise}[Reading a PMF]
\label{ch06:ex-read-pmf}
A random variable $X$ has PMF given by $p_X(0) = 0.1$, $p_X(1) = 0.3$,
$p_X(2) = c$, $p_X(3) = 0.2$, and $p_X(x) = 0$ otherwise.
(a) Find $c$. (b) Compute $\P(X \ge 2)$. (c) Compute the CDF value $F_X(2)$.
(d) Compute $\P(1 \le X < 3)$ and $\P(1 < X \le 3)$, and explain why they differ.
\end{exercise}

\begin{exercise}[Building a random variable]
\label{ch06:ex-build-rv}
Three fair coins are flipped. Let $H$ be the number of heads and let
$L$ be the length of the longest run of consecutive heads (so $L = 0$ if no heads,
$L=3$ for $HHH$). (a) Write out the eight equally likely outcomes and the values of
$H$ and $L$ on each. (b) Give the PMF of $L$. (c) Are $H$ and $L$ the same random
variable? Explain what it means for two random variables on the same $\Omega$ to be
different even when they share a sample space.
\end{exercise}

\begin{exercise}[Max of two dice]
\label{ch06:ex-max-dice}
Roll two fair dice and let $M = \max(a,b)$. (a) Show $\P(M \le m) = (m/6)^2$ for
$m = 1,\dots,6$ by a direct counting argument, and use it to find the CDF.
(b) Obtain the PMF as the jumps of the CDF, i.e.\ $p_M(m) = F_M(m) - F_M(m-1)$, and
verify the masses sum to $1$. (c) Which is more likely, $M = 6$ or $M \le 3$?
\end{exercise}

\begin{exercise}[Indicator algebra]
\label{ch06:ex-indicator-algebra}
Let $A, B, C$ be events. (a) Using only the indicator dictionary
(Proposition~\ref{ch06:prop-indicator-dict}), expand $\ind_{A \cup B \cup C}$ in terms
of $\ind_A, \ind_B, \ind_C$ and their products. (b) Read off the inclusion--exclusion
formula for $\P(A \cup B \cup C)$ by taking the expectation of each side (you may
assume, as we will prove in Chapter~8, that the expectation of a sum of indicators is
the sum of their probabilities). (c) Explain in one sentence why this derivation is
cleaner than a Venn-diagram argument.
\end{exercise}

\begin{exercise}[Transforming a distribution]
\label{ch06:ex-transform}
Let $X$ be uniform on $\set{-3,-2,-1,0,1,2,3}$ (seven values, each mass $1/7$).
(a) Find the PMF of $Y = \abs{X}$. (b) Find the PMF of $Z = X^2 - 2X$. (Watch for
values of $x$ that collide.) (c) Find the PMF of $W = 2X + 1$ and explain why no masses
merge here.
\end{exercise}

\begin{exercise}[Counting fixed points -- a preview]
\label{ch06:ex-fixed-points}
Three people check their coats and the coats are returned in a uniformly random order
(one of the $3! = 6$ permutations). Let $N$ be the number of people who get their own
coat back. (a) For $i = 1,2,3$, let $A_i$ be the event ``person $i$ gets their own
coat''; write $N = \ind_{A_1} + \ind_{A_2} + \ind_{A_3}$ and find $\P(A_i)$ for each.
(b) By listing all six permutations, find the full PMF of $N$. (Notice $N = 2$ is
impossible -- why?) (c) Verify your PMF sums to $1$.
\end{exercise}

\begin{exercise}[CDF properties, conceptual]
\label{ch06:ex-cdf-concept}
Decide whether each function $F$ can be the CDF of some random variable; if not, say
which property of Proposition~\ref{ch06:prop-cdf} fails.
(a) $F(x) = 0$ for $x<0$, $F(x) = x$ for $0 \le x \le 1$, $F(x) = 1$ for $x > 1$.
(b) $F(x) = 0$ for $x < 0$, $F(x) = 1 - e^{-x}$ for $x \ge 0$.
(c) The function that is $0$ for $x<0$, jumps to $0.6$ at $x=0$, drops to $0.4$ at
$x=1$, and is $1$ for $x \ge 2$.
(d) $F(x) = \tfrac{1}{2} + \tfrac{1}{\pi}\arctan(x)$ for all $x \in \R$.
\end{exercise}

\begin{exercise}[Same distribution, different functions]
\label{ch06:ex-same-dist-diff-fn}
On the sample space of a single fair die ($\Omega = \set{1,2,3,4,5,6}$), define
$X(\omega) = \omega$ and $Y(\omega) = 7 - \omega$. (a) Show $X$ and $Y$ have the
\emph{same} distribution. (b) Show that nevertheless $X$ and $Y$ are different
functions, e.g.\ they never agree on any outcome. (c) Compute the PMF of $X + Y$ and
explain why it is degenerate (puts all its mass on one value). What does this teach you
about the difference between ``same distribution'' and ``same random variable''?
\end{exercise}

\begin{exercise}[A proof: distributions determine each other]
\label{ch06:ex-proof-cdf-pmf}
Let $X$ be a discrete random variable taking values in
$\set{x_1 < x_2 < x_3 < \cdots}$. (a) Prove that the PMF determines the CDF via
$F_X(x) = \sum_{x_k \le x} p_X(x_k)$. (b) Prove the converse: the CDF determines the
PMF via $p_X(x_k) = F_X(x_k) - F_X(x_{k}^-)$, where $F_X(x_k^-) = \lim_{t \uparrow x_k}
F_X(t)$ is the left limit. (c) Conclude that two discrete random variables with the
same CDF have the same PMF, and hence the same distribution. Where in your argument did
you use right-continuity?
\end{exercise}

\begin{exercise}[Modeling open-ended]
\label{ch06:ex-modeling}
You roll a fair die repeatedly until the first time a $6$ appears; let $T$ be the
number of rolls (including the final $6$). (a) Argue that the sample space here is
naturally \emph{countably infinite}, and describe it. (b) Find $\P(T = k)$ for
$k = 1, 2, 3, \dots$ and verify $\sum_{k=1}^\infty \P(T=k) = 1$ using the geometric
series. (c) This is your first encounter with a random variable on an infinite sample
space. Explain why the ``masses sum to one'' requirement now genuinely needs the
\emph{infinite} sum, and why that is exactly what countable additivity (Chapter~3) was
built to handle.
\end{exercise}
